\section{Related Work}
\label{sec:related_work}

\para{Self-Reflection and Critique in LLMs.} Iterative self-reflection has been shown to improve performance in many reasoning tasks. \textsc{Reflexion} \citep{shinn2023reflexion} uses verbal reinforcement to help agents learn from past failures, while \textsc{Self-Refine} \citep{madaan2023selfrefine} enables models to iteratively improve their outputs based on their own critiques. While these studies focus on behavioral improvements, our work focuses on the \textit{internal representational shifts} that underlie these changes.

\para{Mechanistic Interpretability and Subspace Structuring.} Recent work has successfully identified linear features and directions in LLM representations that correlate with high-level properties. \citet{tamoyan2025factual} found linear features in the residual stream that indicate whether a model "knows" a fact, while \citet{ackerman2024inspection} identified "self-authorship" vectors that allow models to distinguish their own text. Our work builds on these findings by examining how these representations \textit{drift} dynamically when the model is asked to critique its own internal state.

\para{Representation Engineering and Reflection Control.} The field of Representation Engineering (\repe) \citep{yan2025reflctrl} has proposed methods to directly steer LLM behavior by manipulating activations along specific directions. \citet{zhu2025emergencecontrol} identified "self-reflection vectors" that can be used to modulate the model's reflective tendencies. Our work complements these efforts by analyzing whether natural self-critique (without intervention) induces similar structured shifts in the reasoning subspace, or if it requires more potent external guidance to do so.

\para{Representational Drift.} Drift in LLM representations is often studied in the context of model degradation or fine-tuning \citep{xing2025brainrot}. In this work, we repurpose "drift" as a positive metric for measuring cognitive change during reflection, distinguishing between "shallow" drift (random wandering) and "structured" drift (movement toward a more separable reasoning manifold).
